\chapter{引言}\label{chap:introduction}

不可压缩流体的运动由不可压缩 Navier--Stokes 方程组(incompressible Navier--Stokes equations)描述\index{buke-yasuo-liu@不可压缩流}.该方程组表达了线动量守恒\index{xian-dongliang@线动量守恒}和质量守恒.在稳态\index{wending@稳态流动}且流体黏度较大的特殊情形下,Navier--Stokes 方程组\index{navier-stokes@纳维--斯托克斯方程组}可以用 Stokes 方程(Stokes equations)\index{stokes@斯托克斯方程}近似.不可压缩流问题不仅本身值得研究,也是许多复杂模型的一部分;这些模型用于描述自然现象以及工程和工业过程.

通常无法求得 Stokes 方程或 Navier--Stokes 方程的解析解(analytical solution)\index{jiexi-jie@解析解},因此必须采用数值方法(numerical methods)来近似其解.为此,需要对方程进行离散(discretization)\index{lisan@离散};一般而言,既要作时间离散(temporal discretization),也要作空间离散(spatial discretization).本书讨论空间离散中的有限元方法(finite element method, FEM)\index{youxianyuan-fangfa@有限元方法}.有限元方法应用广泛,从数学角度看,对它的理解也相当深入.

早在 20 世纪 70 年代,人们就开始使用有限元方法模拟 Stokes 方程和 Navier--Stokes 方程.同一时期,对这两类方程的有限元分析也开始发展.例如,\textcite{Babuska1971} 和 \textcite{Brezzi1974} 引入了 inf-sup 条件(inf-sup condition)\index{inf-sup@inf-sup 条件};这一条件是连续问题和有限元问题适定性(well-posedness)的基础.早期有限元分析成果汇集于 \textcite{GiraultRaviart1979} 的专著中.该书的扩充版 \textcite{GiraultRaviart1986} 成为这一领域的经典参考著作.自其出版以来,三十年间出现了许多新的进展,也取得了大量新成果.较近出版、研究不可压缩流问题有限元方法或其重要方面的专著包括 \textcite{Layton2008}、\textcite{Boffi2008} 和 \textcite{Elman2014}.

本书一方面覆盖范围广泛,从 Navier--Stokes 方程的推导一直延伸到湍流模拟;另一方面,有不少主题若要详细阐述,足以自成一部专著,因此本书只作概述.本书着重讨论数学问题.除了有限元方法的许多结果,本书也详细介绍了连续方程的一些基本结果,因为理解连续问题的分析,是全面理解所研究问题的重要基础.

本书的一项主要特点,是对数学工具和大多数证明作了细致讲解.这种写法源于我过去的体会:有时,为了理解文献中常见的简略证明步骤,会花费(甚至浪费)很多时间.其实,若能多给一点提示,或在估计中再多写一行,许多步骤就很容易理解.因此,本书尽可能做到自成体系,使读者理解其中大多数数学结果时无需另查其他文献.总的来说,本书面向广泛的读者群体:本领域的资深研究人员、青年研究人员和硕士生.硕士生也能有效利用本书这一点,已经在教学实践中得到验证.本书的大部分内容曾在柏林自由大学 2013 至 2015 年间开设的硕士课程中讲授;此后,有多篇硕士论文围绕这些课程涉及的主题完成.

\section{本书内容}\label{sec:contents-of-monograph}
\markright{\thesection\ 本书内容}

第 2 章依据质量守恒和线动量守恒,概述不可压缩 Navier--Stokes 方程的推导.随后从物理考虑出发,推导应力张量的重要性质;介绍无量纲方程,并讨论恰当的边界条件.

本书以下各章的安排,围绕第 2 章指出的不可压缩 Navier--Stokes 方程固有的困难展开:
\begin{itemize}
  \item 首先研究速度与压力之间的耦合:
    \begin{itemize}
      \item 第 3 章介绍抽象理论,并讨论如何选取恰当的有限元空间;
      \item 第 4 章将该抽象理论用于 Stokes 方程.
    \end{itemize}
  \item 其次考虑对流占主导这一问题:第 5 章针对 Oseen 方程研究这一主题.Oseen 方程是一类线性化的 Navier--Stokes 方程.
  \item 最后,在前述两项困难之外进一步考虑 Navier--Stokes 方程的非线性:
    \begin{itemize}
      \item 第 6 章研究低雷诺数下出现的稳态流动;
      \item 第 7 章研究中等雷诺数下出现的层流;
      \item 第 8 章研究高雷诺数下出现的湍流.
    \end{itemize}
\end{itemize}

在不可压缩流问题中,速度与压力的耦合意味着不能直接任意选取有限元空间对.为了使问题适定,这些空间必须满足所谓的离散 inf-sup 条件.第 3 章将推导这一条件.推导以抽象线性鞍点问题的适定性研究为基础.首先把抽象理论用于连续线性不可压缩流问题,由此确定速度和压力所适用的函数空间;这些函数空间满足所谓的 inf-sup 条件.接着,本章讨论:连续 inf-sup 条件成立,并不自动保证离散 inf-sup 条件成立.书中给出了不满足该离散条件的速度与压力有限元空间对的例子,随后介绍若干验证离散 inf-sup 条件的技巧,并给出重要的 inf-sup 稳定有限元空间对.对于其中一些空间对,本章还会证明其满足离散 inf-sup 条件,并介绍如何计算离散 inf-sup 常数.本章最后一节讨论 Helmholtz 分解.

第 4 章把上一章建立的理论用于 Stokes 方程.Stokes 方程是线性方程组,也是最简单的不可压缩流模型,只描述由黏性力引起的流动.本章先讨论弱解的存在性、唯一性和稳定性,继而介绍有限元误差分析的结果.误差分析的第一部分研究协调有限元方法,第二部分研究低阶非协调有限元离散;随后给出有限元方法实现方面的一些说明.接下来,本章简要介绍后验误差估计及其在自适应网格加密中的应用.之后讨论绕开离散 inf-sup 条件的若干方法.这些方法使速度和压力可以采用相同的分片多项式空间,从实践角度看颇具吸引力.本章详细分析其中一种方法——PSPG 方法,并简要讨论其他几种方法.一般而言,有限元方法只能近似满足质量守恒.本章最后综述能够减小质量守恒偏差、甚至保证质量守恒的若干方法.

第 5 章研究 Oseen 方程.它是包含黏性(二阶)、对流(一阶)和反应(零阶)项的线性方程.这类方程出现在求解 Navier--Stokes 方程的多种数值方法中,通常对流项或反应项会压过黏性项.分析中的一个主要问题,是追踪稳定性界和误差界对问题系数的依赖.本章先建立方程解的存在性与唯一性,再研究标准 Galerkin 有限元方法.分析表明,当对流或反应占主导时,稳定性界和误差界会变大,数值研究也支持这一结论.为改善数值解,必须采用稳定化方法.本章详细分析一种基于残量的稳定化方法,即 SUPG/PSPG/grad-div 方法,并简要回顾其他若干稳定化方法.

第 6 章研究不可压缩流问题的第一个非线性模型——稳态 Navier--Stokes 方程.首先考察非线性项,介绍其不同形式并推导若干性质.随后研究稳态 Navier--Stokes 方程连续问题的解.结果表明,只有外力足够小且不随时间变化、黏度足够大时,才可期待解唯一存在.针对这一情形,本章给出有限元误差分析,重点在于估计非线性项;然后讨论求解非线性问题的迭代方法.本章最后介绍对偶加权残量(Dual Weighted Residual, DWR)方法,这是一种针对感兴趣的物理量进行后验误差估计的方法.

第 7 章开始研究含时不可压缩 Navier--Stokes 方程.从有限元离散的角度,本章考虑所谓层流,也就是适用标准 Galerkin 有限元方法的流动.开篇简要介绍含时不可压缩 Navier--Stokes 方程弱解的存在性与唯一性分析,并特别阐明为何在实际相关的三维情形下无法证明唯一性.随后介绍连续时间情形下 Galerkin 有限元方法的误差分析;这里的“连续时间”是指尚未对时间作离散.实际模拟必须离散时间,因此本章接下来介绍若干时间步进格式;这些格式在每个离散时刻都要求求解速度与压力耦合问题,其中重点讨论 \(\theta\)-格式.之后以向后 Euler 格式为例,介绍完全离散方程的有限元误差分析.至此介绍的方法都需要求解鞍点问题,这可能带来较高的计算成本.本章最后介绍投影方法,它可以避开鞍点问题:只需分别求解速度场每个分量和压力对应的标量方程.

第 8 章讨论湍流模拟.湍流并没有严格的数学定义,因此本章先描述通常被视为湍流的流场所具有的特征,并勾勒一种描述湍流的数学思路.湍流包含的尺度往往过小,无法在计算成本可接受的网格上分辨;因此必须用所谓湍流模型刻画这些尺度对可分辨尺度的影响.本章主要介绍那些能够进行数学或数值分析,或其推导基于数学论证的湍流模型.大涡模拟(large eddy simulation, LES)是一种非常流行的湍流建模方法,它通过空间平均,只模拟较大的(可解析)尺度.关于 LES 的第一节推导这些尺度对应的方程,尤其考察其中关于微分运算与空间平均可交换的假设;本章指出,通常会出现不可忽略的交换误差.下一节介绍最流行的 LES 模型——Smagorinsky 模型;针对这一模型,已有较完善的数学和数值分析.之后介绍基于波数空间近似推导的 LES 模型,最后讨论近似解卷积模型(Approximate Deconvolution Models, ADM).接着介绍 Leray--\(\alpha\) 模型,该模型通过正则化速度场构造.随后讨论 Navier--Stokes--\(\alpha\) 模型,其推导基于对拉格朗日泛函及相应轨线的研究.本章还介绍变分多尺度(Variational Multiscale, VMS)方法.与 LES 不同,VMS 方法通过适当函数空间中的投影来定义需要模拟的大尺度.VMS 方法主要分为两类:基于两尺度分解的方法,以及基于三尺度分解的方法.本章详细介绍五种不同的 VMS 方法实现,最后给出若干湍流数值模拟研究,其中使用了 Smagorinsky 模型和一种有代表性的 VMS 模型.

对不可压缩 Navier--Stokes 方程进行线性化和离散后,许多方法会得到速度与压力耦合的代数方程组.第 9 章简要介绍这类方程组的求解器,包括稀疏直接求解器和迭代求解器;后者需要配合恰当的预条件子.本章着重介绍一种耦合多重网格方法,它被用于本书大多数数值算例.

附录 A 汇集泛函分析的一些基本记号、若干不等式和定理;这些结果用于本书的分析和数值分析.附录 B 介绍有限元方法的一些基础知识,特别详细描述用于离散不可压缩流问题的有限元空间.附录 C 讨论如何通过插值或投影逼近 Sobolev 空间中的函数,并给出有限元误差分析中大量使用的相应估计.附录 D 介绍若干数值模拟算例,并将它们分为三组:稳态流问题算例、含时层流问题算例和湍流问题算例.这些算例用于开展本书所呈现的数值模拟.原书目录还列有附录 E“记号”,本中文版将依原书一并收录.

柏林自由大学开设的硕士课程覆盖以下主题:
\begin{itemize}
  \item 课程 1:第 2、3 章,以及第 4 章第 4.1--4.3 节;
  \item 课程 2:第 4 章第 4.4--4.6 节,第 5--7 章,以及第 9 章;
  \item 课程 3:第 8 章.
\end{itemize}
当然,在这些课程中,讲授内容只集中于各主题中最重要的问题.
